\chapter{Heart Failure}

\section*{Overview}
Heart failure is a chronic condition where the heart doesn't pump blood as well as it should. Despite the name, it doesn't mean the heart has stopped working. Rather, the heart's pumping action is weakened, causing blood to back up in the veins and organs. Over time, the heart muscle weakens and stretches, becoming less efficient at pumping blood throughout the body.

\section*{Symptoms}
Common symptoms of heart failure include:
\begin{itemize}
    \item Shortness of breath (dyspnea), especially with activity or when lying flat
    \item Fatigue and weakness
    \item Swelling (edema) in legs, ankles, and feet
    \item Rapid or irregular heartbeat
    \item Reduced exercise tolerance
    \item Persistent cough or wheezing with white or pink blood-tinged mucus
    \item Difficulty breathing when lying flat
    \item Sudden weight gain from fluid retention
    \item Impaired thinking or confusion
    \item Lack of appetite and nausea
\end{itemize}
Symptoms may be subtle initially and worsen gradually.


\section*{Causes}
Heart failure commonly results from:
\begin{itemize}
    \item \textbf{Coronary artery disease:} Most common cause, leading to heart attack
    \item \textbf{High blood pressure:} Forces heart to work harder
    \item \textbf{Cardiomyopathy:} Damage to heart muscle from various causes
    \item \textbf{Heart valve disease:} Valves don't work properly
    \item \textbf{Congenital heart defects:} Present at birth
    \item \textbf{Arrhythmias:} Abnormal heart rhythms
    \item \textbf{Diabetes:} Increases cardiovascular risk
    \item \textbf{Viral infections:} Can damage heart muscle
    \item \textbf{Alcohol and drug abuse:} Toxic to heart muscle
\end{itemize}

\section*{Diagnosis}
Heart failure diagnosis involves:
\begin{itemize}
    \item \textbf{Physical examination:} Checking for swelling, lung sounds, heart murmurs
    \item \textbf{Blood tests:} BNP or NT-proBNP (heart failure markers), complete blood count, kidney function
    \item \textbf{Echocardiogram:} Gold standard - measures ejection fraction (EF)
    \item \textbf{Chest X-ray:} Shows heart size and fluid in lungs
    \item \textbf{Electrocardiogram:} Detects arrhythmias, previous heart attacks
    \item \textbf{Cardiac MRI:} Detailed heart structure and function
    \item \textbf{Stress test:} Evaluates heart function during exercise
    \item \textbf{Coronary angiography:} Checks for blocked arteries
\end{itemize}
Heart failure with reduced EF (HFrEF) vs preserved EF (HFpEF) guides treatment.


\section*{Treatment}
Heart failure treatment includes:
\begin{itemize}
    \item \textbf{Medications:}
        \begin{itemize}
            \item ACE inhibitors or ARBs (lisinopril, losartan)
            \item Beta-blockers (metoprolol, carvedilol)
            \item Diuretics (furosemide) for fluid removal
            \item Aldosterone antagonists (spironolactone)
            \item SGLT2 inhibitors (dapagliflozin, empagliflozin)
            \item ARNI (sacubitril/valsartan) for eligible patients
        \end{itemize}
    \item \textbf{Devices:} Pacemakers, ICDs, biventricular pacemakers
    \item \textbf{Surgery:} Coronary bypass, valve repair/replacement, heart transplant
    \item \textbf{Lifestyle changes:} Salt restriction, fluid management, exercise
    \item \textbf{Cardiac rehabilitation:} Supervised exercise and education
\end{itemize}

\section*{Prevention}
Prevent heart failure by:
\begin{itemize}
    \item Managing high blood pressure
    \item Controlling diabetes
    \item Keeping cholesterol in check
    \item Maintaining healthy weight
    \item Exercising regularly
    \item Eating a balanced diet
    \item Avoiding smoking and excessive alcohol
    \item Getting adequate sleep
    \item Managing stress
    \item Regular health checkups
\end{itemize}
Early detection and management of cardiovascular disease reduces heart failure risk.


\clearpage